\heading{实验物理中的统计方法-第1章 基本概念}


\begin{question}
考虑某样本空间 \(S\) 以及给定子空间 \(B\)，并假设 \(P(B)>0\)。证明条件概率

\[
P(A|B)=\frac{P(A\cap B)}{P(B)}.
\]

满足概率的公理。
\begin{solution}
对固定的 \(B\) 且 \(P(B)>0\)，定义
\(P_B(A)=P(A|B)=P(A\cap B)/P(B)\)。验证概率公理：

\begin{enumerate}
\item
  非负性：\(P(A\cap B)\ge0\)，故 \(P_B(A)\ge0\)。
\item
  归一性：\(P_B(S)=P(S\cap B)/P(B)=P(B)/P(B)=1\)。
\item
  可列可加性：若 \(A_i\) 两两不交，则 \((A_i\cap B)\) 也两两不交，故
\end{enumerate}

\[
P_B\left(\bigcup_i A_i\right)=\frac{P\left(\bigcup_i(A_i\cap B)\right)}{P(B)}
=\frac{\sum_iP(A_i\cap B)}{P(B)}=\sum_iP_B(A_i).
\]

因此条件概率满足概率公理。
\end{solution}
\end{question}


\begin{question}
证明

\[
P(A\cup B)=P(A)+P(B)-P(A\cap B).
\]

(提示：将 \(A\cup B\) 表示成 3 个不相交的子集的并。)
\begin{solution}
把 \(A\cup B\) 分成三个互不相交的部分：

\[
A\cup B=(A\setminus B)\cup(B\setminus A)\cup(A\cap B).
\]

于是

\[
P(A\cup B)=P(A\setminus B)+P(B\setminus A)+P(A\cap B).
\]

又

\[
P(A)=P(A\setminus B)+P(A\cap B),\qquad
P(B)=P(B\setminus A)+P(A\cap B),
\]

所以

\[
P(A\cup B)=P(A)+P(B)-P(A\cap B).
\]
\end{solution}
\end{question}


\begin{question}
某粒子束流包含 \(10^{-4}\) 的电子，其余为光子。粒子通过某双层探测器，可能在两层都给出信号，也可能只有一层给出信号或者没有任何信号。电子 \((e)\) 和光子 \((\gamma)\) 在穿过该双层探测器给出 0，1 或 2 个信号的概率如下

\[
\begin{split}
P(0|e)&=0.001,\qquad P(0|\gamma)=0.99899,\\
P(1|e)&=0.01,\qquad P(1|\gamma)=0.001,\\
P(2|e)&=0.989,\qquad P(2|\gamma)=10^{-5}.
\end{split}
\]

\begin{enumerate}
\item
  如果只有一层给出信号，该粒子为光子的概率是多少？
\item
  如果两层都给出了信号，该粒子为电子的概率是多少？
\end{enumerate}
\begin{solution}
设先验概率 \(P(e)=10^{-4}\)，\(P(\gamma)=1-10^{-4}=0.9999\)。

\begin{enumerate}
\item
  只有一层给出信号时，用 Bayes 公式：
\end{enumerate}

\[
P(\gamma|1)=\frac{P(1|\gamma)P(\gamma)}{P(1|e)P(e)+P(1|\gamma)P(\gamma)}.
\]

代入数值：

\[
P(\gamma|1)=\frac{0.001\times0.9999}{0.01\times10^{-4}+0.001\times0.9999}
\approx0.9990.
\]

\begin{enumerate}[start=2]
\item
  两层都给出信号时：
\end{enumerate}

\[
P(e|2)=\frac{P(2|e)P(e)}{P(2|e)P(e)+P(2|\gamma)P(\gamma)}
=\frac{0.989\times10^{-4}}{0.989\times10^{-4}+10^{-5}\times0.9999}
\approx0.9082.
\]
\end{solution}
\end{question}


\begin{question}
假设随机变量 \(x\) 的概率密度函数为 \(f(x)\)。证明 \(y=x^2\) 的概率密度函数为

\[
g(y)=\frac{1}{2\sqrt{y}}f(\sqrt{y})+\frac{1}{2\sqrt{y}}f(-\sqrt{y}).
\]
\begin{solution}
令 \(y=x^2\)。当 \(y>0\) 时，反解有两个分支 \(x=\sqrt y\) 和
\(x=-\sqrt y\)，且

\[
\left|\frac{dx}{dy}\right|=\frac{1}{2\sqrt y}.
\]

变量变换公式给出

\[
g(y)=f(\sqrt y)\frac{1}{2\sqrt y}+f(-\sqrt y)\frac{1}{2\sqrt y},\qquad y>0.
\]

即

\[
g(y)=\frac{1}{2\sqrt y}f(\sqrt y)+\frac{1}{2\sqrt y}f(-\sqrt y).
\]
\end{solution}
\end{question}


\begin{question}
假设两个独立的随机变量 \(x\) 和 \(y\) 都服从 0 到 1 之间的均匀分布，即概率密度函数 \(g(x)\) 为

\[
g(x)=\begin{cases}
1, & 0<x<1,\\
0, & \text{其它}.
\end{cases}
\]

概率密度函数 \(h(y)\) 与 \(g(x)\) 类似。

\begin{enumerate}
\item 利用 Statistical Data Analysis 中的 (1.35) 式，证明 \(z=xy\) 的概率密度函数 \(f(z)\) 为
\[
f(z)=\begin{cases}
-\log z, & 0<z<1,\\
0, & \text{其它}.
\end{cases}
\]

\begin{enumerate}[start=2]
\item
  利用 Statistical Data Analysis 的 (1.37) 和 (1.38)
  式，通过另外定义一个函数 \(u=x\)，求 \(z=xy\) 的概率密度函数。首先求 \(z\) 和
  \(u\) 的联合概率密度函数，然后对 \(u\) 进行积分求出 \(z\) 的概率密度函数。
\item
  证明 \(z\) 的累积分布为
\[
  F(z)=z(1-\log z).
  \]
\end{enumerate}
\end{enumerate}
\begin{solution}
\begin{enumerate}
\item
  对 \(z=xy\)，且 \(x,y\sim U(0,1)\) 独立。固定 \(x\)，由
  \(y=z/x\)，需要 \(0<z/x<1\)，即 \(z<x<1\)。因此
\end{enumerate}

\[
f_Z(z)=\int_z^1 1\cdot1\cdot\frac{1}{x}\,dx=-\log z,\qquad 0<z<1.
\]

其它 \(z\) 区间密度为 0。

\begin{enumerate}[start=2]
\item
  令 \(u=x\)，\(z=xy\)，则 \(x=u\)，\(y=z/u\)。Jacobian 为
\end{enumerate}

\[
\left|\frac{\partial(x,y)}{\partial(z,u)}\right|=
\begin{vmatrix}
0&1\\
1/u&-z/u^2
\end{vmatrix}=\frac{1}{u}.
\]

由于 \(0<u<1\) 且 \(0<z/u<1\)，有 \(z<u<1\)。所以联合密度

\[
f_{Z,U}(z,u)=\frac1u,\qquad 0<z<u<1.
\]

积分掉 \(u\)：

\[
f_Z(z)=\int_z^1\frac1u\,du=-\log z.
\]

\begin{enumerate}[start=3]
\item
  累积分布为
\end{enumerate}

\[
F(z)=P(XY\le z)=\int_0^z[-\log t]dt=z(1-\log z),\qquad 0<z<1.
\]
\end{solution}
\end{question}


\begin{question}
考虑随机变量 \(x\) 与常数 \(\alpha\) 和 \(\beta\)。证明

\[
\begin{split}
E[\alpha x+\beta]&=\alpha E[x]+\beta,\\
V[\alpha x+\beta]&=\alpha^2 V[x].
\end{split}
\]
\begin{solution}
利用期待值线性性：

\[
E[\alpha x+\beta]=\alpha E[x]+\beta.
\]

方差为

\[
V[\alpha x+\beta]=E[(\alpha x+\beta-\alpha E[x]-\beta)^2]
=E[\alpha^2(x-E[x])^2]=\alpha^2V[x].
\]
\end{solution}
\end{question}


\begin{question}
考虑两个随机变量 \(x\) 和 \(y\)。

\begin{enumerate}
\item
  证明 \(\alpha x+y\) 的方差为
\end{enumerate}

\[
\begin{split}
V[\alpha x+y]&=\alpha^2 V[x]+V[y]+2\alpha\operatorname{cov}[x,y]\\
&=\alpha^2 V[x]+V[y]+2\alpha\rho\sigma_x\sigma_y.
\end{split}
\]

其中 \(\alpha\) 为任意常数，\(\sigma_x^2=V[x]\)，\(\sigma_y^2=V[y]\)，关联系数 \(\rho=\operatorname{cov}[x,y]/(\sigma_x\sigma_y)\)。

\begin{enumerate}[start=2]
\item 利用 (a) 的结果，证明关联系数总是位于区间 \(-1\le \rho\le 1\)。(利用 \(V[\alpha x+y]\) 的方差总是大于或等于零。)
\end{enumerate}
\begin{solution}
\begin{enumerate}
\item
  由方差定义：
\end{enumerate}

\[
\begin{split}
V[\alpha x+y]
&=E[(\alpha(x-E[x])+(y-E[y]))^2]\\
&=\alpha^2V[x]+V[y]+2\alpha\operatorname{cov}[x,y]\\
&=\alpha^2\sigma_x^2+\sigma_y^2+2\alpha\rho\sigma_x\sigma_y.
\end{split}
\]

\begin{enumerate}[start=2]
\item
  因为方差非负，对任意 \(\alpha\)，二次函数
\end{enumerate}

\[
\alpha^2\sigma_x^2+2\alpha\rho\sigma_x\sigma_y+
\sigma_y^2\ge0.
\]

其判别式必须不大于 0：

\[
(2\rho\sigma_x\sigma_y)^2-4\sigma_x^2\sigma_y^2\le0.
\]

若 \(\sigma_x\sigma_y>0\)，得到 \(\rho^2\le1\)，即
\(-1\le\rho\le1\)。若某个方差为零，相关系数退化，结论可由极限理解。
\end{solution}
\end{question}


\begin{question}
假设随机变量 \(x=(x_1,\ldots,x_n)^T\) 用联合概率密度函数 \(f(x)\) 描述，而变量 \(y=(y_1,\ldots,y_n)^T\) 由下面的线性变换定义

\[
y_i=\sum_{j=1}^{n} A_{ij}x_j.
\]

假设反变换 \(x=A^{-1}y\) 存在。

\begin{enumerate}
\item 证明 \(y\) 的联合概率密度函数为
\[
g(y)=f(A^{-1}y)\left|\det(A^{-1})\right|.
\]

\begin{enumerate}[start=2]
\item
  当 \(A\) 为正交矩阵，即 \(A^{-1}=A^T\) 时，求 \(g(y)\)。
\end{enumerate}
\end{enumerate}
\begin{solution}
\begin{enumerate}
\item
  线性变换为 \(y=Ax\)，反变换 \(x=A^{-1}y\)。多维变量变换公式给出
\end{enumerate}

\[
g(y)=f(A^{-1}y)\left|\det\frac{\partial x}{\partial y}\right|
=f(A^{-1}y)|\det(A^{-1})|.
\]

\begin{enumerate}[start=2]
\item
  若 \(A\) 为正交矩阵，则 \(A^{-1}=A^T\)，且 \(|\det(A^{-1})|=1\)。因此
\end{enumerate}

\[
g(y)=f(A^Ty).
\]
\end{solution}
\end{question}

